\documentclass[11pt]{article}
\usepackage[margin=1in]{geometry}
\usepackage{amsmath,amssymb}
\usepackage[hidelinks]{hyperref}
\setlength{\parindent}{0pt}
\setlength{\parskip}{0.7em}
\begin{document}

\begin{center}
{\Large\bfseries Literature resolution of the Wang--Madiman Rényi EPI conjecture}\par
\vspace{0.4em}
{\small Source: arXiv:1604.04225 \qquad Answer: arXiv:1904.08038}
\end{center}

\textbf{Status.} Literature already answered; no new mathematical result is
claimed here.

\textbf{Source question.} In Section 2.3.3, pp.~14--16, of Madiman,
Melbourne, and Xu, \emph{Forward and Reverse Entropy Power Inequalities in
Convex Geometry}, the Wang--Madiman conjecture is stated as follows.  Let
$X_1,\ldots,X_n$ be independent random vectors in $\mathbb R^d$, let
$p>d/(d+2)$, and choose independent scaled generalized Gaussians $Z_i$ so
that $h_p(X_i)=h_p(Z_i)$.  Then
\[
 N_p(X_1+\cdots+X_n)\geq N_p(Z_1+\cdots+Z_n).
\]
The survey reports that all cases other than the classical $p=1$ case and
the established $p=\infty$ cases remained open.

\textbf{Exact later answer.} Jaye, Livshyts, Paouris, and Pivovarov,
\emph{Remarks on the Rényi Entropy of a sum of IID random variables},
arXiv:1904.08038, restate this assertion as Conjecture~1.1.  Their
Introduction states that it fails already when
\[
 d=1,\qquad p=2,\qquad n=2,
\]
with $X_1,X_2$ independent and identically distributed.  Section~4 applies
the Euler--Lagrange equation derived in Section~3 and shows that the proposed
generalized Gaussian does not satisfy the necessary first-order condition.
Consequently it is not a minimizer of the Rényi entropy of the sum under the
fixed-entropy constraint, disproving the conjecture.

\textbf{Scope.} The conjectured extremizer is decisively ruled out, but the
answering paper leaves the identity of the true minimizer open.

\end{document}

